\section{Distributed sources and their forward operator}\label{sec:sources}

\subsection{Source distributions precede an array}
Fix a regular interface, material state and frequency satisfying the harmonic
assumptions of \cref{sec:acoustics}. Let $\Gamma_a\subseteq\partial\mathcal D_c$
be the permitted boundary source support, and let $D_s\subset\mathcal D_c$ be
the permitted volume source support. Either can have several components and
can be empty. Define the complex peak source fields
\begin{equation}
 u=(g,\bm F,q),\qquad
 [g]=\mathrm{m\,s^{-1}},\quad [\bm F]=\mathrm{N\,m^{-3}},
 \quad [q]=\mathrm{s^{-1}}.
 \label{eq:distributed-control}
\end{equation}
Here $g$ is imposed outward boundary velocity, $\bm F$ is a first-order
body-force density, and $q$ is an optional volumetric dilation rate in the
linear acoustic equations. Volume fields are understood phase by phase and
vanish outside $D_s$. No restriction to radial, annular, or equal-phase
sources is made. Their physical time signals are the real parts of their
phasors multiplied by $e^{-\mathrm i\omega t}$.

The dilation source is a generalized acoustic representation, not an extra
unexplained mass supply in \cref{eq:mass}. An implementation must derive it
from its mechanism, such as a moving source volume or heat input, and supply
the associated mean mass, energy and thermal balances. In a closed,
mass-conserving mechanical realization one can set $q=0$ and use boundary
motion and body forcing. An acoustic source description alone does not specify
its mean mechanical reaction or heat production.

For definiteness choose positive reference area $S_a$, source volume $V_s$,
and amplitudes $g_0,F_0,q_0$. Define a complex Hilbert source space
$\mathcal U$ with squared norm
\begin{equation}
 \norm{u}_{\mathcal U}^2=
 \frac{1}{S_ag_0^2}\int_{\Gamma_a}|g|^2\dd A+
 \frac{1}{V_sF_0^2}\int_{D_s}|\bm F|^2\dd V+
 \frac{1}{V_sq_0^2}\int_{D_s}|q|^2\dd V.
 \label{eq:source-norm}
\end{equation}
Omit components that are not available. This norm is dimensionless; it is not
automatically power. The corresponding complex inner product is conjugate
linear in its first argument. Spatial smoothness can instead be penalized
with a Sobolev norm. The admissible subset $\mathcal U_{\rm ad}$ specifies
support, amplitude, smoothness and any hardware restrictions. Temporal
bandwidth is a condition on a waveform or envelope, not on a single phasor.

The displayed weak and trace equations assume regular enough source fields
and solutions. In the explicit derivation, $\bm F,q$ are supported away from
the interface and outer wall. Separated source and observation sets permit
additional interior regularity; singular point sources need a different
function-space treatment. Measure-valued Helmholtz source inversion is one
such treatment \cite{pieper2020}; its existence and sparse-recovery results
do not automatically apply to a quadratic interface-stress objective.

\subsection{An explicit transmission problem}
The following is a useful, limited member of the general model: constant
real $\rho_\ell,c_\ell$ in each phase, negligible first-order capillary
compliance, negligible mean advection, and inviscid bulk acoustics. Dissipation
is represented here by a declared boundary admittance. The forced equations are
\begin{equation}
 -\mathrm i\omega\rho_\ell\bm V_\ell=-\nabla P_\ell+\bm F_\ell,
 \qquad
 -\frac{\mathrm i\omega}{\rho_\ell c_\ell^2}P_\ell
       =-\nabla\cdot\bm V_\ell+q_\ell.
 \label{eq:forced-first-order}
\end{equation}
Define $\alpha_\ell=1/\rho_\ell$ and
$\beta_\ell=\omega^2/(\rho_\ell c_\ell^2)$. Eliminating velocity gives
\begin{equation}
 L_\ell P_\ell:=-\nabla\cdot(\alpha_\ell\nabla P_\ell)-\beta_\ell P_\ell
       =-\nabla\cdot(\alpha_\ell\bm F_\ell)-\mathrm i\omega q_\ell.
 \label{eq:forced-helmholtz}
\end{equation}
The colon-equals sign defines the operator $L_\ell$. With $\partial_n$ taken
along the same minus-to-plus normal on both sides, transmission is
\begin{equation}
 \jump{P}=0,\qquad \jump{\alpha\partial_nP}=0,
 \qquad v=\frac{\alpha_\ell\partial_nP_\ell}{\mathrm i\omega}.
 \label{eq:transmission-source}
\end{equation}
The scalar $v$ is the common acoustic normal-velocity trace; its units are
velocity. Pressure continuity here is the stated high-frequency reduction,
not a replacement for the full interface conditions at all frequencies.

Let $\bm n_a$ be the domain-outward normal on the cylinder boundary, and let
$Y(\bm x,\omega)$ be a prescribed acoustic admittance with units
$\mathrm{m\,Pa^{-1}\,s^{-1}}$. Set $g=0$ outside $\Gamma_a$ and impose
\begin{equation}
 \bm V\cdot\bm n_a=YP+g,
 \qquad
 \alpha\partial_{n_a}P=\mathrm i\omega(YP+g).
 \label{eq:source-robin}
\end{equation}
The second equation uses the stated absence of $\bm F$ near the wall.
The case $Y=0$ is prescribed wall velocity. A passive local admittance has
$\Rea Y\geq0$; arbitrary elastic wall dynamics need a coupled operator in
place of this scalar relation. An open exterior also needs its own radiation
condition. Even with passive boundaries, well-posedness at the chosen frequency
must be established; an undamped cavity at an eigenfrequency is not covered
by the inverse operator used below.

For a complex, globally continuous test pressure $\psi\in H^1(\mathcal D_c)$,
where $H^1$ means a square-integrable function with square-integrable first
derivatives, integration by parts yields $a_\Gamma(P,\psi)=b(u,\psi)$, with
\begin{align}
 a_\Gamma(P,\psi)&=
  \sum_\ell\int_{\Omega^\ell}
   (\alpha_\ell\nabla P_\ell\cdot\nabla\bar\psi
                         -\beta_\ell P_\ell\bar\psi)\dd V
      -\mathrm i\omega\int_{\partial\mathcal D_c}YP\bar\psi\dd A,
      \label{eq:acoustic-weak-a}\\
 b(u,\psi)&=\mathrm i\omega\int_{\Gamma_a}g\bar\psi\dd A
       +\int_{D_s}\alpha\bm F\cdot\nabla\bar\psi\dd V
       -\mathrm i\omega\int_{D_s}q\bar\psi\dd V.
       \label{eq:acoustic-weak-b}
\end{align}
A bar denotes complex conjugation. The weak form defines exactly how the
distributed controls enter the wave operator. It can be replaced systematically
by a thermoviscous or solid-coupled weak form; those replacements also change
the load and its adjoint. Established active Helmholtz control illustrates the
importance of source support, resonance and regularization
\cite{hubenthal2016}. That literature controls compatible wave fields in stated
regions; it is not a theorem of arbitrary fluid-shape controllability.

\subsection{Input power is an independent consistency condition}
Define the net time-averaged acoustic source power entering this canonical
problem by
\begin{equation}
 P_{\rm src}=-\frac12\Rea\int_{\Gamma_a}P\bar g\dd A
    +\frac12\Rea\int_{D_s}\bm F\cdot\bar{\bm V}\dd V
    +\frac12\Rea\int_{D_s}P\bar q\dd V.
 \label{eq:source-power}
\end{equation}
The minus sign for $g$ follows from its outward convention: positive work into
the fluid uses traction $-P\bm n_a$. Every term has units of watts. In a settled
state with lossless bulk coefficients, the power identity is
\begin{equation}
 \boxed{P_{\rm src}=
       \frac12\int_{\partial\mathcal D_c}\Rea Y\,|P|^2\dd A.}
 \label{eq:power-identity}
\end{equation}
To derive it, put $\psi=P$ in \cref{eq:acoustic-weak-a,eq:acoustic-weak-b},
take $-\Ima/(2\omega)$, and substitute
$\bm V=\alpha(\nabla P-\bm F)/(\mathrm i\omega)$.
The term $\bm F\cdot\bar{\bm F}$ contributes no real source work in that
substitution. Real bulk loss, mean work, or unresolved wall layers require
their corresponding terms in a more complete power balance. Net acoustic
power is distinct from gross electrical consumption and from
\cref{eq:source-norm}.

\subsection{Cylinder symmetry does not impose source symmetry}
Let $(r,\theta,z)$ be cylindrical coordinates, with $r^2=x_1^2+x_2^2$.
Only when geometry, material state and boundary operators are invariant under
rotation can the linear wave problem be separated into integer azimuthal
orders $m\in\mathbb Z$, proportional to $e^{\mathrm i m\theta}$. Products
of an order-$m$ field with the conjugate of an order-$n$ field have order
$m-n$. Thus quadratic radiation loads contain differences of the excited
angular orders. One pure angular order produces an axisymmetric mean normal
load in this invariant configuration; coherent combinations can produce
non-axisymmetric loads. Axisymmetric source fields alone cannot produce a
non-axisymmetric load in this fixed, axisymmetric linear wave problem.
Once the interface or boundary breaks that symmetry, angular orders couple.
This selection rule follows from the displayed product, and is a testable
consequence of the source formulation rather than an array-design assumption.
